\section{Random variables and discrete distributions}
\subsection*{Core problems}
\begin{problembook}
A player stakes 5 leva and rolls two independent standard dice. Two sixes pay
100 leva; exactly one six pays 5 leva; no sixes pay nothing. Let $X$ be the player's
net gain. Find the distribution and expectation of $X$. Is the game fair?
\end{problembook}
\begin{problembook}
A person tosses a fair coin, taking a step forwards for heads and backwards for
tails. Toss outcomes are independent. What is the probability that after 10 tosses the person is:
\begin{enumeratea}
  \item at the starting point;
  \item at distance 2 from it;
  \item five steps ahead of it?
\end{enumeratea}
\end{problembook}
\begin{problembook}
Let $Z_1$ and $Z_2$ be the outcomes of two independent dice rolls, and let $X=Z_1+Z_2$.
\begin{enumeratea}
  \item The dice are standard. Find the distribution, expectation and variance of $X$.
  \item Let
  \[
  \Pbb(Z_i=1)=\Pbb(Z_i=6)=\frac14
  \quad\text{and}\quad\Pbb(Z_i=j)=\frac18,
  \qquad\text{for }i\in\{1,2\},\ j\in\{2,3,4,5\}.
  \]
  Find the distribution, expectation and variance of $X$, showing two ways of finding the expectation.
  \item For standard dice, let $Y=|X-7|$. Find the distribution of $Y$ and $\Ebb Y$.
  Compare $\Ebb Y$ with $|\Ebb X-7|$, and $\Ebb Y^2$ with $\Var X$.
\end{enumeratea}
\end{problembook}
\begin{problem}
A box contains 7 light bulbs, 3 of them defective. Four distinct bulbs are chosen
uniformly. Find the distribution and expectation of the number of defective bulbs
chosen. Also find the expectation using indicators, without summing over the whole
distribution. Find the variance using indicators as well. Are these indicators independent?
\end{problem}
\begin{problembook}
Two independent standard dice are rolled five times, with outcomes in different
rolls also independent. What is the probability that exactly 2 rolls have sum 6?
Find the expectation and variance of the number of such rolls.
\end{problembook}
\begin{problempencil}[Memorylessness; Test 1, SE 2023]
Let $X\sim\Geom(p)$, meaning $\Pbb(X=k)=p(1-p)^k$ for $k\in\{0,1,2,\dots\}$,
where $p\in(0,1]$.
\begin{enumeratea}
  \item Prove that for all integers $m,n\ge0$ with $\Pbb(X\ge n)>0$,
  \[ \Pbb(X\ge m+n\mid X\ge n)=\Pbb(X\ge m). \]
  \item Conversely, let $Y$ take values in $\{0,1,2,\dots\}$ and satisfy the above
  equality for every pair of integers $m,n\ge0$ with $\Pbb(Y\ge n)>0$.
  Prove that $Y\sim\Geom(p)$ for some $p\in(0,1]$.
\end{enumeratea}
\end{problempencil}
\begin{problembook}
Assume that the number of minor earthquakes in a region over $t$ months has a
Poisson distribution with parameter $2t$. What is the probability of fewer than
four earthquakes in three months?
\end{problembook}
\begin{probleminteresting}[St. Petersburg paradox]
A casino charges an entry fee of $A$ leva. We toss a fair coin until the first
heads, with independent outcomes. If heads first occurs on toss $n$, the prize
is $2^n$ leva. Find the expected prize. For which $A$ would you personally play?
What changes if the prize is $\min\{2^n,2^{30}\}$ leva? How does this result
compare with your answer about the values of $A$ for which you would play?
\end{probleminteresting}
\subsection*{Additional problems}
\begin{problem}
An urn contains 5 white and 3 black balls. We draw until a white ball appears.
Let $X$ be the number of black balls drawn. Find its distribution, expectation and variance for:
\begin{enumeratea}
  \item sampling without replacement;
  \item sampling with replacement.
\end{enumeratea}
\end{problem}
\begin{problem}
$80\%$ of home printers work without additional configuration, with independent
outcomes for different printers and weeks. A company sells 10 printers each week.
Find the probability that at least 9 work immediately; the probability that this
happens in five consecutive weeks; and the probability that the first week in
which fewer than 9 work immediately is week $k$. Calculate the last probability
for $k=3$ and $k=21$, and find the expected number of weeks until the first such week.
\end{problem}
\begin{problempencil}[Test 1, SE 2023]
\begin{enumeratea}
  \item Let $X_1,\dots,X_n$ be independent and identically distributed, with
  $F(m)=\Pbb(X_1\le m)$, and let $X_{\min}=\min\{X_1,\dots,X_n\}$ and
  $X_{\max}=\max\{X_1,\dots,X_n\}$. Express $\Pbb(X_{\min}\le m)$ and
  $\Pbb(X_{\max}\le m)$ in terms of $F(m)$ and $n$.
  \item Five standard dice are rolled. In each round, we roll all remaining dice
  simultaneously and remove those showing six. All rolls are independent. Let $K$
  be the number of rounds until all dice have been removed. What is the distribution
  of the number of rolls of one die? Find $\Pbb(K\le n)$ and $\Pbb(K=n)$ for every integer $n\ge1$.
\end{enumeratea}
\end{problempencil}
\begin{problem}
Let $X$ take values in $\{0,1,2,\dots\}$ and have finite expectation.
\begin{enumeratea}
  \item Prove that
  \[ \Ebb X=\sum_{k=1}^\infty\Pbb(X\ge k)=\sum_{n=0}^\infty\Pbb(X>n). \]
  \item Check the formula for $X\sim\Geom(p)$.
  \item Find $\Ebb K$ for the number of rounds $K$ in the preceding problem and
  compare it with the expected number of rolls of one die.
\end{enumeratea}
\end{problem}
\begin{problem}
Independent Bernoulli trials have success probability $p$, where $0<p<1$.
Find the probability that the $r$th success occurs on exactly trial $k+r$, for
integers $r\ge1$ and $k\ge0$. Let $X_i$ indicate success on trial $i$. Express
the event using $X_i$ and the partial sums $S_n=\sum_{i=1}^nX_i$, with $S_0=0$.
\end{problem}
\begin{problem}
A shot hits its target with probability $0.001$. There are 5000 independent shots,
and at least two hits are required. Find the exact probability and the Poisson
approximation. How large is the numerical difference?
\end{problem}
\begin{problem}
A submarine fires $n$ times; shots are independent and each torpedo hits with
probability $p$. On a hit, one of the $m$ compartments is selected uniformly,
independently of other hits. Find the probability that at least two distinct compartments are flooded.
\end{problem}
\begin{problem}[Banach's matchboxes]
A smoker carries two boxes containing $n$ matches each. Each time, the smoker
independently chooses one box with probability $1/2$ and takes a match.
On first reaching for an empty box, what is the probability that the other box
contains exactly $k\in\{0,\dots,n\}$ matches?
\end{problem}
\subsection*{Challenges and independent research}
\begin{problemstar}
Let $\pi$ be a uniformly random permutation of $\{1,\dots,n\}$, and let
\[ I_n=\#\{(i,j):1\le i<j\le n,\ \pi(i)>\pi(j)\}. \]
Prove that
\[ \Ebb I_n=\frac{n(n-1)}4\qquad\text{and}\qquad\Var(I_n)=\frac{n(n-1)(2n+5)}{72}. \]
\end{problemstar}
\begin{problemstar}[Coupon collector]
Let $X_1,X_2,\dots$ be independent and identically distributed, uniform on
$\{1,\dots,m\}$, where $m\ge2$, and let
\[ T=\inf\{n\ge1:\{X_1,\dots,X_n\}=\{1,\dots,m\}\}. \]
\begin{enumeratea}
  \item Find $\Pbb(T\le n)$ for $n\ge1$.
  \item Prove that
  \[ \Ebb T=mH_m,\qquad\Var(T)=m^2H_m^{(2)}-mH_m, \]
  where $H_m=\sum_{j=1}^m\frac1j$ and $H_m^{(2)}=\sum_{j=1}^m\frac1{j^2}$.
\end{enumeratea}
Where do you use independence of the waiting times, and where is linearity of expectation alone sufficient?
\end{problemstar}
\begin{researchtask}[Generalised coupon collector]
Let $X_1,X_2,\dots$ be independent and identically distributed random variables
taking values in $[m]=\{1,\dots,m\}$, where $m\ge2$, with
\[
\Pbb(X_n=i)=p_i>0,\qquad\text{for }i\in[m],
\qquad\text{and }\sum_{i=1}^m p_i=1.
\]
Let $T$ be the number of draws until all $m$ values have appeared.
\begin{enumeratea}
  \item Prove that
  \[
  \Ebb T=\sum_{\emptyset\ne A\subseteq[m]}\frac{(-1)^{|A|+1}}{\sum_{i\in A}p_i}
  \qquad\text{and}\qquad
  \Ebb T=\int_0^\infty\left(1-\prod_{i=1}^m\bigl(1-e^{-p_it}\bigr)\right)\,dt.
  \]
  Explain why the first formula is unsuitable for direct numerical calculation
  when $m=30$. Suggest a numerical method for the integral and a way to check its accuracy.
  \item Compare values from the integral formula with simulations for $m=10,30,100$
  in the following models:
  \begin{enumerate}[label=(\roman*)]
    \item uniform probabilities: $p_i=1/m$ for $i\in[m]$;
    \item probabilities
    \[ p_i=\frac{i^{-\alpha}}{\sum_{j=1}^m j^{-\alpha}},\qquad\text{for }i\in[m], \]
    considering both $\alpha=1$ and $\alpha=2$;
    \item one rare coupon:
    \[ p_1=m^{-2}\quad\text{and}\quad p_i=\frac{1-m^{-2}}{m-1},\qquad\text{for }i\in\{2,\dots,m\}. \]
  \end{enumerate}
  State the number of simulations and discuss the accuracy of the estimated mean.
  Explain which mechanism determines the scale of $\Ebb T$ in each model.
\end{enumeratea}
You may use AI for the proof and numerical method, provided you check the results
and explain how you checked them.
\end{researchtask}
\noindent\textbf{Lean.} Let $S$ be a finite nonempty set, $X,Y:S\to\mathbb R$,
and $a,b\in\mathbb R$. On the finite uniform probability space on $S$, define
\[ \Ebb X=\frac1{|S|}\sum_{s\in S}X(s). \]
Formalise in Lean 4 with Mathlib and prove linearity:
\[ \Ebb(aX+bY)=a\,\Ebb X+b\,\Ebb Y. \]
First prove, by induction on $S$ and without using \texttt{Finset.sum\_add\_distrib}, that
\[ \sum_{s\in S}\bigl(X(s)+Y(s)\bigr)=\sum_{s\in S}X(s)+\sum_{s\in S}Y(s); \]
\texttt{Finset.mul\_sum} and \texttt{Finset.sum\_div} are useful for the rest.
Why is nonemptiness of $S$ needed to define $\Ebb$, but not for the identity for sums?
